
\documentclass[12pt,a4paper]{article}
\usepackage[utf8]{inputenc}
\usepackage[T5]{fontenc}
\usepackage[vietnamese]{babel}
\usepackage{amsmath}
\usepackage{graphicx}
\usepackage{listings}
\usepackage{xcolor}
\usepackage{longtable}
\usepackage{geometry}
\usepackage{hyperref}
\usepackage{tikz}
\usetikzlibrary{positioning,arrows.meta}
\usepackage{booktabs}
\usepackage{float}
\usepackage{fancyhdr}
\usepackage{lastpage}
\usepackage{setspace}
\usepackage{titlesec}
\geometry{a4paper,left=2cm,right=2cm,top=2.5cm,bottom=2.5cm}

\definecolor{codebg}{rgb}{0.95,0.95,0.92}
\definecolor{hcmutblue}{RGB}{3,37,126}

\lstset{
  backgroundcolor=\color{codebg},
  basicstyle=\ttfamily\small,
  breaklines=true,
  frame=single,
  numbers=left,
  numberstyle=\tiny\color{gray},
  keywordstyle=\color{blue}\bfseries,
  commentstyle=\color{green!60!black}\itshape,
  stringstyle=\color{red},
  showstringspaces=false,
  tabsize=2
}

\pagestyle{fancy}
\fancyhead[L]{\small Trường ĐH Bách Khoa Tp.HCM --- Khoa KHKT Máy Tính}
\fancyfoot[L]{\small Computer Networks - Lab Report}
\fancyfoot[R]{\small Trang \thepage/\pageref{LastPage}}
\setstretch{1.25}
\titleformat{\section}{\Large\bfseries\color{hcmutblue}}{\thesection}{1em}{}
\titleformat{\subsection}{\large\bfseries}{\thesubsection}{1em}{}

\begin{document}

%===== BÌA =====
\begin{titlepage}
\begin{center}
  \vspace*{20pt}
  \large\textbf{ĐẠI HỌC QUỐC GIA THÀNH PHỐ HỒ CHÍ MINH}\\[5pt]
  \Large\textbf{TRƯỜNG ĐẠI HỌC BÁCH KHOA}\\[30pt]
  \textcolor{hcmutblue}{\fontsize{22pt}{22pt}\selectfont\textbf{COMPUTER NETWORKS}}\\[10pt]
  \fontsize{16pt}{16pt}\selectfont\textbf{Bài tập lớn 1}\\[40pt]
\end{center}
\begin{table}[H]
\centering\large
\renewcommand{\arraystretch}{1.5}
\begin{tabular}{>{\bfseries}ll}
  Instructor: & Bùi Xuân Giang \\
  Class ID:   & CO3094 (A01) \\
  Students:   &
    \begin{tabular}[t]{@{}l@{\hspace{1cm}}l@{}}
      Nguyễn Chí Thanh  & 2313078 \\
      Lê Đức Tài        & 2312995 \\
      Nguyễn Việt Hoàng & 2311066 \\
      Mai Xuân Nhựt     & 2312549 \\
    \end{tabular}
\end{tabular}
\end{table}
\vfill
\begin{center}\large\textbf{HỒ CHÍ MINH, 2026}\end{center}
\end{titlepage}

\newpage\tableofcontents\newpage

%===== 1. TỔNG QUAN =====
\section{Tổng quan kiến trúc hệ thống}

Hệ thống được xây dựng theo mô hình đa tiến trình gồm ba thành phần chính:
\begin{itemize}
  \item \textbf{Backend Server} (\texttt{daemon/backend.py}): Xử lý HTTP request, xác thực, phục vụ file tĩnh và API RESTful.
  \item \textbf{ChatApp / AsynapRous} (\texttt{apps/chatapp.py}, \texttt{daemon/asynaprous.py}): Ứng dụng chat kết hợp Client--Server và Peer-to-Peer.
  \item \textbf{Tracker} (\texttt{apps/trackerapp.py}): Quản lý đăng ký và khám phá peer.
\end{itemize}

Toàn bộ hệ thống sử dụng \textbf{Python standard library} (\texttt{socket}, \texttt{select}, \texttt{threading}) mà \textbf{không dùng bất kỳ web framework hay asyncio bên ngoài nào}. Thay vào đó, nhóm tự xây dựng một \textbf{Event Loop} dựa trên \texttt{select()} hệ thống.

%===== 2. NON-BLOCKING =====
\section{Cơ chế Non-blocking (Phần 2.1)}

\subsection{Hai chế độ xử lý đồng thời}

Hệ thống hỗ trợ hai cơ chế được lựa chọn qua tham số \texttt{--mode}:

\begin{longtable}{|p{4cm}|p{6.5cm}|p{4.5cm}|}
\hline
\textbf{Cơ chế} & \textbf{Mô tả} & \textbf{Lệnh} \\ \hline
\textbf{Threading} & Mỗi kết nối tạo một daemon thread riêng & \texttt{--mode threading} \\ \hline
\textbf{Callback} (mặc định) & Custom \texttt{select()} Event Loop, 1 thread phục vụ nhiều kết nối & \texttt{--mode callback} \\ \hline
\end{longtable}

\subsection{Chế độ Callback --- Custom \texttt{select()} Event Loop}

Đây là đóng góp trọng tâm của bài tập: tự xây dựng Event Loop sử dụng \textbf{lời gọi hệ thống \texttt{select.select()}} thay vì dùng \texttt{asyncio}.

\subsubsection*{Kiến trúc Event Loop (\texttt{daemon/eventloop.py})}

\begin{lstlisting}[language=Python, caption={EventLoop --- vòng lặp sự kiện tự xây}]
class EventLoop:
    """
    Vong lap su kien (singleton) dua tren select.select().
    1 thread duy nhat phuc vu nhieu ket noi dong thoi.
    """
    def _run_once(self):
        # Buoc 1: Chay pending callbacks (TaskQueue)
        self._task_queue.drain()
        # Buoc 2: Hoi OS socket nao san sang (non-blocking)
        r_list = self._registry.get_read_sockets()
        readable, writable, _ = select.select(
            r_list, [], [], self._select_timeout
        )
        # Buoc 3: Goi callback tuong ung
        for sock in readable:
            cb = self._registry.get_read_callback(sock)
            if cb:
                cb(sock)

    def run_forever(self):
        self._running = True
        while self._running:
            self._run_once()
\end{lstlisting}

\textbf{Nguyên lý 1 iteration (1 tick):}
\begin{enumerate}
  \item Drain \texttt{TaskQueue} --- chạy các callback đã lên lịch.
  \item Gọi \texttt{select.select()} --- hỏi OS danh sách socket sẵn sàng I/O.
  \item Gọi callback đã đăng ký cho từng socket sẵn sàng.
  \item Lặp lại --- không bao giờ dừng cho đến khi \texttt{stop()}.
\end{enumerate}

\subsubsection*{CallbackRegistry --- Danh bạ socket → callback}

\begin{lstlisting}[language=Python]
class CallbackRegistry:
    """Anh xa: socket -> callback function."""
    def register_read(self, sock, callback):
        self._read_callbacks[sock] = callback
    def unregister(self, sock):
        self._read_callbacks.pop(sock, None)
\end{lstlisting}

\subsubsection*{SelectHTTPServer --- HTTP Server tích hợp EventLoop}

\begin{lstlisting}[language=Python, caption={Luong xu ly ket noi moi}]
class SelectHTTPServer:
    def start(self):
        self._server_sock = socket.socket(AF_INET, SOCK_STREAM)
        self._server_sock.setblocking(False)   # Non-blocking!
        self._server_sock.bind((self.ip, self.port))
        self._server_sock.listen(128)
        self.loop.register_read(self._server_sock, self._on_new_connection)

    def _on_new_connection(self, server_sock):
        conn, addr = server_sock.accept()
        conn.setblocking(False)
        self._buffers[conn] = ConnectionBuffer(conn, addr)
        self.loop.register_read(conn, self._on_data)

    def _on_data(self, conn):
        data = conn.recv(4096)   # Non-blocking
        buf = self._buffers[conn]
        buf.feed(data)
        if buf.done:             # Da nhan du 1 HTTP request
            self._handle_request(conn, buf)

    def _handle_request(self, conn, buf):
        adapter = HttpAdapter(self.ip, self.port, None, buf.addr, self.routes)
        response = adapter.process_request(buf.get_request(), buf.addr)
        conn.sendall(response)
        self._close_conn(conn)
\end{lstlisting}

\textbf{Phân tầng trách nhiệm:}
\begin{itemize}
  \item \textbf{EventLoop}: Quản lý vòng lặp \texttt{select()}, phân phối sự kiện.
  \item \textbf{SelectHTTPServer}: Quản lý kết nối TCP (accept, recv, send).
  \item \textbf{ConnectionBuffer}: Gom dữ liệu từng mảnh đến khi đủ 1 HTTP request.
  \item \textbf{HttpAdapter.process\_request()}: Xử lý HTTP (parse, auth, routing).
\end{itemize}

\subsubsection*{Chế độ Threading}

\begin{lstlisting}[language=Python]
# daemon/backend.py --- threading mode
if mode_async == "threading":
    client_thread = threading.Thread(
        target=handle_client,
        args=(ip, port, conn, addr, routes)
    )
    client_thread.daemon = True
    client_thread.start()
\end{lstlisting}

Mỗi kết nối được giao cho một daemon thread xử lý độc lập. Phù hợp khi số kết nối đồng thời thấp, mỗi request cần nhiều CPU.

\subsection{Khởi động server}

\begin{lstlisting}[language=bash]
# Callback mode (mac dinh) --- port 9020
python start_backend.py --server-port 9020 --mode callback

# Threading mode --- port 9010
python start_backend.py --server-port 9010 --mode threading
\end{lstlisting}

%===== 3. XÁC THỰC =====
\section{Xác thực HTTP (Phần 2.2)}

\subsection{Luồng xác thực trong \texttt{HttpAdapter.process\_request()}}

Toàn bộ logic HTTP (parse + auth + routing) được tập trung trong method \texttt{process\_request()}, được gọi bởi cả 3 mode:

\begin{lstlisting}[language=Python, caption={process\_request() --- loi HTTP da dung chung}]
def process_request(self, raw_msg, addr):
    req.prepare(raw_msg, self.routes)   # Parse HTTP

    # Buoc 1: Kiem tra Cookie session_id (RFC 6265)
    if req.cookies and 'session_id' in req.cookies:
        sid = req.cookies['session_id']
        if sid in ACTIVE_SESSIONS:
            is_authenticated = True
            current_user = ACTIVE_SESSIONS[sid]
        else:
            user = validate_session(sid)
            if user:
                is_authenticated = True
                add_session(sid, user)

    # Buoc 2: HTTP Basic Auth (RFC 7235)
    if not is_authenticated and req.auth:
        decoded = base64.b64decode(auth_parts[1]).decode()
        username, password = decoded.split(':', 1)
        if VALID_USERS.get(username) == password:
            new_session = create_session(username)
            add_session(new_session, username)
            new_cookie_to_set = f"session_id={new_session}; Path=/; HttpOnly"

    # Buoc 3: Phan quyen
    if not is_authenticated and not is_public:
        return b"HTTP/1.1 401 Unauthorized\r\n..."  # hoac 302 redirect

    # Buoc 4: Routing -> API handler
    response = master_api_handler(req, resp)
    return response
\end{lstlisting}

\subsection{Quản lý Session}

\begin{longtable}{|p{4.5cm}|p{10.5cm}|}
\hline
\textbf{Thành phần} & \textbf{Mô tả} \\ \hline
\texttt{ACTIVE\_SESSIONS} & Dictionary RAM: \texttt{session\_id $\rightarrow$ username} \\ \hline
\texttt{db/sessions\_id.txt} & Lưu persistent: \texttt{sid|user|expire\_timestamp} \\ \hline
\texttt{db/account.txt} & Tài khoản: \texttt{username:password} \\ \hline
\texttt{SESSION\_TTL} & 86400 giây (24 giờ) \\ \hline
\end{longtable}

\begin{longtable}{|p{4cm}|p{2cm}|p{9cm}|}
\hline
\textbf{Endpoint} & \textbf{Method} & \textbf{Mô tả} \\ \hline
\texttt{/api/login} & POST & Xác thực, trả \texttt{Set-Cookie: session\_id=...} \\ \hline
\texttt{/api/logout} & POST & Xóa session, đặt \texttt{Max-Age=0} \\ \hline
\texttt{/api/me} & GET & Trả tên user hiện tại \\ \hline
\end{longtable}

%===== 4. HTTP =====
\section{Xử lý HTTP Request và Response}

\subsection{Module \texttt{Request} (\texttt{daemon/request.py})}

\begin{lstlisting}[language=Python]
def prepare(self, request, routes=None):
    self.method, self.path, self.version = self.extract_request_line(request)
    self.headers = self.prepare_headers(request)
    self.body    = self.fetch_headers_body(request)[1]
    # Parse Cookie
    self.cookies = {}
    for pair in cookie_str.split(';'):
        if '=' in pair:
            k, v = pair.strip().split('=', 1)
            self.cookies[k] = v
    self.auth = self.headers.get('authorization', None)
    if routes:
        self.hook = routes.get((self.method, self.path))
\end{lstlisting}

\subsection{Module \texttt{HttpAdapter} (\texttt{daemon/httpadapter.py})}

Sau khi refactor, \texttt{HttpAdapter} có 3 method chính:
\begin{itemize}
  \item \texttt{process\_request(raw\_msg, addr)}: Lõi HTTP --- dùng chung cho cả 3 mode.
  \item \texttt{handle\_client(conn, addr, routes)}: Threading mode --- recv() rồi gọi \texttt{process\_request()}.
  \item \texttt{handle\_client\_coroutine(reader, writer)}: Async I/O --- \texttt{await read()} rồi gọi \texttt{process\_request()}.
\end{itemize}

%===== 5. CHATAPP =====
\section{Ứng dụng Chat Hybrid (Phần 2.3)}

\subsection{Kiến trúc Hybrid Client-Server + P2P}

\begin{center}
\begin{tikzpicture}[node distance=1.5cm,
  box/.style={draw,rectangle,align=center,minimum width=2.8cm,minimum height=1.2cm,font=\ttfamily\small},
  arrow/.style={->,>=Stealth,thick}]
  \node[box](tracker){Tracker\\(port 9001)};
  \node[box,below left=1.5cm and 0.5cm of tracker](peerA){Peer A\\(port 8001)};
  \node[box,below=2cm of tracker](peerB){Peer B\\(port 8002)};
  \node[box,below right=1.5cm and 0.5cm of tracker](peerC){Peer C\\(port 8003)};
  \draw[arrow](tracker)--(peerA) node[midway,left,font=\footnotesize]{register};
  \draw[arrow](tracker)--(peerB);
  \draw[arrow](tracker)--(peerC);
  \draw[<->,thick](peerA.south)--+(0,-0.4)-|(peerC.south)
    node[midway,below,font=\footnotesize\ttfamily]{Direct P2P (send-peer)};
\end{tikzpicture}
\end{center}

\subsection{Giai đoạn khởi tạo (Client--Server)}

\begin{lstlisting}[language=Python, caption={trackerapp.py --- dang ky peer}]
@app.route('/submit-info', methods=['POST'])
def submit_info(req):
    resolved_ip = resolve_peer_ip(ip, req.remote_addr)
    PEER_REGISTRY[name] = {"ip": resolved_ip, "port": port, "last_seen": now}

@app.route('/get-list', methods=['POST'])
def get_list(req):
    cleanup_online(now)   # Xoa peer offline (TTL 30s)
    peers = [{"name": n, **info} for n, info in PEER_REGISTRY.items()
             if ONLINE.get(n)]
    return json_response({"code": 1, "peers": peers})
\end{lstlisting}

\subsection{Giai đoạn Chat (Peer-to-Peer, không dùng asyncio)}

ChatApp mới dùng \textbf{blocking socket} thay vì \texttt{asyncio.open\_connection()}, và dùng \texttt{threading} thay vì \texttt{asyncio.gather()}:

\begin{lstlisting}[language=Python, caption={send\_to\_peer\_sync() --- gui P2P bang blocking socket}]
def send_to_peer_sync(target_ip, target_port, payload):
    body = json.dumps(payload).encode('utf-8')
    http_request = (
        f"POST /internal/receive-msg HTTP/1.1\r\n"
        f"Content-Length: {len(body)}\r\nConnection: close\r\n\r\n"
    ).encode() + body
    s = socket.socket(AF_INET, SOCK_STREAM)
    s.settimeout(5)
    s.connect((target_ip, int(target_port)))
    s.sendall(http_request)
    s.recv(4096)    # Doc va bo qua response
    s.close()
\end{lstlisting}

\begin{lstlisting}[language=Python, caption={broadcast\_peer() --- gui dong thoi bang threading}]
@app.route('/broadcast-peer', methods=['POST'])
def broadcast_peer(req):
    payload = {"from": sender, "message": message, ...}
    threads = []
    for name, info in PEER_CACHE.items():
        if name == sender_name: continue
        t = threading.Thread(
            target=send_to_peer_sync,
            args=(info["ip"], info["port"], payload),
            daemon=True
        )
        t.start(); threads.append(t)
    for t in threads: t.join(timeout=5)
\end{lstlisting}

\begin{lstlisting}[language=Python, caption={Nhan va lay tin nhan}]
@app.route('/internal/receive-msg', methods=['POST'])
def internal_receive_msg(req):
    MESSAGE_QUEUE.append(json.loads(req.body))
    return json_response({"code": 1, "message": "Received"})

@app.route('/poll-messages', methods=['POST'])
def poll_messages(req):
    msgs = list(MESSAGE_QUEUE)
    MESSAGE_QUEUE.clear()
    return json_response({"code": 1, "messages": msgs})
\end{lstlisting}

%===== 6. BENCHMARK =====
\section{Benchmark: Threading vs Callback}

\subsection{Phương pháp}

Script \texttt{benchmark\_th\_cb.py} gửi $N$ request đến 2 server, đo throughput và latency:

\begin{lstlisting}[language=bash]
python benchmark_th_cb.py --th-port 9010 --cb-port 9020 -n 1000 --step 50 --max-c 500
\end{lstlisting}

\subsection{Kết quả đo lường}

\begin{longtable}{|r|r|r|r|r|l|}
\hline
\textbf{Concurrency} & \textbf{TH RPS} & \textbf{CB RPS} & \textbf{TH Avg(ms)} & \textbf{CB Avg(ms)} & \textbf{Winner} \\ \hline
50  & 1010 & 1547 & 36.8 & 13.5 & Callback \\ \hline
100 & 1049 & 1603 & 72.2 & 25.2 & Callback \\ \hline
150 & 273  & 1556 & 148  & 37.3 & Callback \\ \hline
200 & 344  & 343  & 200  & 136  & Callback \\ \hline
300 & 405  & 501  & 251  & 234  & Callback \\ \hline
400 & 416  & 339  & 314  & 333  & Threading \\ \hline
500 & 443  & 419  & 388  & 437  & Threading \\ \hline
\end{longtable}

\subsection{Phân tích}

\begin{itemize}
  \item \textbf{Callback thắng} khi concurrency thấp--vừa ($\leq$300): 1 thread không tốn overhead tạo thread mới, \texttt{select()} xử lý nhiều kết nối ngắn rất hiệu quả. Latency trung bình thấp hơn 2--4 lần.
  \item \textbf{Threading thắng} khi concurrency rất cao ($>$400): Callback bị bottleneck vì 1 thread phải tuần tự xử lý từng callback. Thread pool cho phép song song thực sự trên multi-core.
  \item \textbf{Callback đặc biệt vượt trội} ở concurrency 150 (1556 vs 273 RPS): Do threading tạo quá nhiều thread đồng thời gây context switch overhead.
\end{itemize}

%===== 7. CẤU TRÚC FILE =====
\section{Cấu trúc file dự án}

\begin{lstlisting}[language=bash]
MMT_Assignment1/
├── daemon/
│   ├── eventloop.py     # Custom select() Event Loop (MỚI)
│   ├── backend.py       # 3 mode: threading / callback
│   ├── httpadapter.py   # process_request() dung chung
│   ├── request.py       # HTTP parser
│   ├── response.py      # HTTP response builder
│   ├── asynaprous.py    # Micro-framework routing (khong asyncio)
│   └── api.py           # API handlers
├── apps/
│   ├── chatapp.py       # P2P chat (sync socket + threading)
│   └── trackerapp.py    # Peer registry
├── www/                 # HTML pages
├── db/                  # account.txt, sessions_id.txt
├── start_backend.py     # Entry: --mode threading/callback
├── start_chatapp.py     # Entry: chat peer
├── start_tracker.py     # Entry: tracker
├── benchmark_th_cb.py   # Benchmark threading vs callback
└── test_chatapp.py      # Test toan bo chuc nang
\end{lstlisting}

%===== 8. KẾT LUẬN =====
\section{Kết luận}

Bài tập lớn đã thực hiện được các yêu cầu chính:

\begin{enumerate}
  \item \textbf{Non-blocking HTTP Server tự xây}: Xây dựng \texttt{EventLoop} dựa trên \texttt{select.select()} trong \texttt{daemon/eventloop.py}, \textbf{không dùng asyncio hay bất kỳ framework bên ngoài nào}. Tuân thủ tinh thần ``develop your own framework using the provided non-blocking mechanisms'' của đề bài.

  \item \textbf{Phân tầng rõ ràng}:
    \begin{itemize}
      \item Tầng I/O: \texttt{EventLoop} + \texttt{SelectHTTPServer} (select).
      \item Tầng HTTP: \texttt{HttpAdapter.process\_request()} (dùng chung 3 mode).
      \item Tầng API: \texttt{master\_api\_handler()} + routes.
    \end{itemize}

  \item \textbf{Xác thực HTTP}: HTTP Basic Auth (RFC 7235) + Cookie session (RFC 6265), phân quyền public/private/API path.

  \item \textbf{Chat Hybrid}: Client--Server (Tracker) + P2P trực tiếp dùng blocking socket + threading cho broadcast đồng thời.

  \item \textbf{Benchmark}: Callback mode vượt trội ở concurrency thấp--vừa (tối đa $\times$5.7 RPS so với threading ở c=150), Threading phù hợp hơn ở concurrency rất cao.

  \item \textbf{AsynapRous Framework}: Micro-framework decorator routing tự implement, không phụ thuộc thư viện ngoài.
\end{enumerate}

\end{document}
